\section{Discussion}
\label{sec:discussion}

\paragraph{ML-DSA as a triage stress test.}
The ML-DSA-65 result (Section~\ref{sec:results-coverage}) is a stress test for
conservative triage, not a leak claim. The debug build names only registered
rejection-sampling functions, but optimized builds inline those functions into
\texttt{crypto\_sign\_signature\_ctx} and also surface \texttt{pack\_sig}. CT-KAT
therefore leaves the row at \emph{needs-analysis}: the registry explains part of
the finding, while the build matrix exposes attribution that is too coarse for
automatic acceptance. This is the intended behavior for a screening artifact: it
turns ambiguous evidence into a precise review item rather than a stronger claim.

\paragraph{Why complementary candidate sources, not one.}
KyberSlash shows that structural CT checking and instruction-latency scanning
are genuinely orthogonal: Memcheck~\cite{valgrind,ctgrind} passes the variant
clean while asm-scan surfaces the secret-derived division candidate~\cite{kyberslash}.
Neither refines the other; each sees inputs the other is blind to. A single
checker, however sound on its own class~\cite{ctverif,binsecrel}, would miss the
rest --- which is why the workflow binds several candidate sources before triage.

\paragraph{Limitations.}
The structural and timing checks are dynamic and see only exercised paths;
untaken-branch leaks are out of scope. For KEMs specifically, the structural
harness decapsulates a \emph{valid} ciphertext, so Valgrind covers the normal
path only; the implicit-rejection (Fujisaki--Okamoto) path---an important ML-KEM
leak surface---is not structurally exercised in this corpus and can only be explored
by dudect's \texttt{leak\_target:~fo} timing mode rather than by the structural
check. The dudect test~\cite{dudect} runs
under QEMU/Docker with a noisy cycle counter (Table~\ref{tab:dudect}); large
separations are unambiguous but precise statistics should be confirmed
natively (\texttt{taskset}, frequency scaling off). Asm-scan is candidate-level
and warn-only: flagging a variable-latency instruction~\cite{kocher} is not
proof its operand is secret, so findings feed triage, not a verdict. The
synthetic harnesses isolate one feature each and are deliberately not PQC, and
the corpus covers one ML-KEM~\cite{fips203,kyber} and one
ML-DSA~\cite{pqclean} family. CT-KAT reflects these scope boundaries in its
reports: clean rows are screening results, not absence proofs, and \emph{ct-leak}
requires an explicit human verdict.
